\documentclass[11pt]{article}
\usepackage[letterpaper,margin=0.88in]{geometry}
\usepackage[T1]{fontenc}
\usepackage{lmodern}
\usepackage{microtype}
\usepackage{amsmath,amssymb,amsthm,mathtools}
\usepackage{graphicx}
\usepackage{xcolor}
\usepackage{enumitem}
\usepackage{booktabs}
\usepackage[hidelinks]{hyperref}
\usepackage{fancyhdr}

\definecolor{navy}{HTML}{17365D}
\definecolor{blue}{HTML}{2A6496}
\definecolor{pale}{HTML}{EEF4F9}
\definecolor{green}{HTML}{236B45}
\hypersetup{colorlinks=true,linkcolor=navy,citecolor=green,urlcolor=blue}
\setlist{nosep,leftmargin=1.4em}
\setlength{\parindent}{0pt}
\setlength{\parskip}{0.52em}
\pagestyle{fancy}
\setlength{\headheight}{14pt}
\fancyhf{}
\lhead{Harmonic Langevin orthogonal dynamics}
\rhead{arXiv:2109.01990 -- candidate full result}
\cfoot{\thepage}
\renewcommand{\headrulewidth}{0.35pt}

\newtheorem{theorem}{Theorem}
\newtheorem{lemma}{Lemma}
\newtheorem{corollary}{Corollary}
\theoremstyle{remark}
\newtheorem{remark}{Remark}

\newcommand{\HH}{\mathsf H}
\newcommand{\K}{\mathcal K}
\newcommand{\Pj}{\mathcal P}
\newcommand{\Qj}{\mathcal Q}
\newcommand{\E}{\mathbb E}
\newcommand{\R}{\mathbb R}

\title{\vspace{-2.4em}\textbf{Exact Gap Doubling for Harmonic Langevin\\Mori--Zwanzig Orthogonal Dynamics}}
\author{Candidate solution packet for Zhu, arXiv:2109.01990\\
\small agent\_lane\_03 \quad Model GPT5.6}
\date{26 August 2026}

\begin{document}
\maketitle
\vspace{-1.1em}

\begin{center}
\fcolorbox{navy}{pale}{\parbox{0.91\textwidth}{
\textbf{Status: candidate full result, likely valid.}
For a natural rank-one Mori projection in a genuinely degenerate harmonic
Langevin system, the optimal positive orthogonal decay gap is exactly twice
the original gap.  The fluctuation force and memory kernel have an even
faster exact exponential law.}}
\end{center}

\section{Source question and convention repair}

Zhu's Section 5, open problem 1, asks whether the Mori--Zwanzig orthogonal
dynamics can rigorously have a larger spectral gap and decay faster for a
degenerate stochastic system \cite{Zhu2021}.  The exact source passage is on
PDF page 10.  The source writes the desired inequality as
$\lambda_{\Qj}<\lambda$, while
its earlier estimates have the form $Ce^{-\lambda t}$.  With positive decay
exponents, faster decay means the reverse inequality.  We therefore make the
convention explicit.  For a semigroup $U(t)$ on an invariant Hilbert space
$E$, set
\[
 \operatorname{gap}(U;E):=
 \sup\{\omega\geq0:\ \exists C<\infty\ \text{with }
 \|U(t)|_E\|\leq C e^{-\omega t}\ \forall t\geq0\}.
\]
Thus a larger positive gap means faster decay.  In the signed spectral-edge
convention, faster decay means that the rightmost nonzero spectral edge moves
to the left.  The theorem below proves both equivalent statements.

\clearpage
\begin{center}
\includegraphics[width=0.96\textwidth]{figures/open_problem_crop.png}

\small\textit{Source evidence: Section 5 and all of open problem 1, including
its convention-revealing footnote, on the original PDF page 10.}
\end{center}

\section{Model and full result}

Fix $\gamma>2$ and consider the one-dimensional harmonic Langevin equation
\begin{equation}\label{eq:sde}
 dq_t=p_t\,dt,\qquad
 dp_t=-q_t\,dt-\gamma p_t\,dt+\sqrt{2\gamma}\,dW_t.
\end{equation}
Its backward generator is
\begin{equation}\label{eq:generator}
 \K=p\partial_q-q\partial_p-\gamma p\partial_p
       +\gamma\partial_p^2
\end{equation}
on $H=L^2(\mu)$, where
$d\mu=(2\pi)^{-1}e^{-(q^2+p^2)/2}\,dq\,dp$.  The diffusion in
\eqref{eq:sde} has rank one in a two-dimensional phase space, so the system
is genuinely degenerate.  Its drift and noise nevertheless satisfy the
Kalman condition.

Let $u(q,p)=q$.  Since $\|q\|_H=1$, the natural scalar Mori projection and
its complement are
\[
 \Pj f=\langle f,q\rangle_H q,\qquad \Qj=I-\Pj.
\]
Write $T(t)=e^{t\K}$ and let $T_{\Qj}(t)$ be the orthogonal semigroup generated
by the closure of $\Qj\K\Qj$ on
$E=\Qj H\cap\{1\}^{\perp}$.  Finally put
\begin{equation}\label{eq:ab}
 a=\frac{\gamma-\sqrt{\gamma^2-4}}2,
 \qquad
 b=\frac{\gamma+\sqrt{\gamma^2-4}}2.
\end{equation}
Then $0<a<b$, $a+b=\gamma$, and $ab=1$.

\begin{theorem}[Exact acceleration]\label{thm:main}
For the degenerate stochastic system \eqref{eq:sde},
\begin{align*}
 \operatorname{gap}(T;\{1\}^{\perp})&=a,\\
 \operatorname{gap}(T_{\Qj};E)&=2a.
\end{align*}
Equivalently, the rightmost nonzero spectral edge moves from $-a$ for
$\K$ to $-2a$ for $\Qj\K\Qj$.  In particular, the positive orthogonal gap is
strictly larger and in fact exactly doubles.

Moreover, the Mori fluctuation force $F(t)$ and scalar memory kernel $M(t)$
for $u=q$ are
\begin{equation}\label{eq:force-memory}
 F(t)=T_{\Qj}(t)\Qj\K q=e^{-\gamma t}p,
 \qquad
 M(t)=\langle q,\K F(t)\rangle_H=-e^{-\gamma t}.
\end{equation}
Thus these two quantities decay at the still faster rate $\gamma>2a$.
\end{theorem}

\section{Proof intuition}

The Gaussian $L^2$ space splits orthogonally into Wiener chaoses.  Harmonic
Langevin evolution preserves every chaos, and the $n$th-chaos semigroup is
the $n$th symmetric tensor power of the two-dimensional first-chaos
semigroup.  Consequently the slowest centered mode occurs in the first
chaos.

The Mori projection onto $q$ deletes one direction from that first chaos.
The remaining first-chaos direction $p$ is driven by the compressed scalar
generator $-\gamma$.  Crucially, the projection is zero on every higher
chaos, so the slowest untouched mode is the tensor square of the original
slow mode.  Its exponent is exactly $2a$.  This explains both the doubling
and its sharpness; no hypocoercive estimate is being compared to another
nonsharp estimate.

\section{Proof}

Let
\[
 B=\begin{pmatrix}0&1\\-1&-\gamma\end{pmatrix},
 \qquad e_p=\begin{pmatrix}0\\1\end{pmatrix}.
\]
Then \eqref{eq:sde} is $dX_t=BX_t\,dt+\sqrt{2\gamma}\,e_p\,dW_t$.
The identity
\[
 B+B^{\mathsf T}+2\gamma e_pe_p^{\mathsf T}=0
\]
proves invariance of the standard Gaussian.  Also $e_p$ and
$Be_p=(1,-\gamma)^{\mathsf T}$ span $\R^2$, proving the Kalman condition.

Let $\HH_n$ be the $n$th Gaussian Wiener chaos, so
$H=\bigoplus_{n\geq0}\HH_n$, $\HH_0=\operatorname{span}\{1\}$, and
$\HH_1=\operatorname{span}\{q,p\}$.  We next derive, rather than assume, the
semigroup action on these spaces.

\begin{lemma}[Exact chaos action]\label{lem:chaos}
Set $S(t)=e^{tB^{\mathsf T}}$ on $\R^2\cong\HH_1$.  Then
\[
 T(t)|_{\HH_n}=S(t)^{\odot n},
 \qquad
 \|T(t)|_{\HH_n}\|=\|S(t)\|^n
 \quad(n\geq1),
\]
where $\odot n$ denotes the $n$th symmetric tensor power.
\end{lemma}

\begin{proof}
For $z\in\R^2$, introduce the Wick exponential
\[
 \mathcal E_z(x)=\exp\!\left(z\cdot x-\frac{|z|^2}{2}\right).
\]
Conditionally on $X_0=x$, $X_t=e^{tB}x+Y_t$, where $Y_t$ is centered Gaussian
with covariance $I-e^{tB}e^{tB^{\mathsf T}}$.  Gaussian integration gives
\[
 T(t)\mathcal E_z(x)
 =\exp\!\left((S(t)z)\cdot x-\frac{|S(t)z|^2}{2}\right)
 =\mathcal E_{S(t)z}(x).
\]
Comparing homogeneous terms in $z$ proves the symmetric-tensor identity.
The upper norm bound is the ordinary tensor-power bound.  Equality follows
by choosing a unit vector at which the two-dimensional operator $S(t)$
attains its norm and taking its symmetric pure tensor.
\end{proof}

Because $T(t)$ is a Markov contraction, $\|S(t)\|\leq1$.  Lemma
\ref{lem:chaos} therefore yields the exact centered norm identity
\begin{equation}\label{eq:centered-norm}
 \|T(t)|_{\{1\}^{\perp}}\|
 =\sup_{n\geq1}\|S(t)\|^n=\|S(t)\|.
\end{equation}
On the ordered basis $(q,p)$, the first-chaos generator is
\begin{equation}\label{eq:first}
 \K|_{\HH_1}=B^{\mathsf T}
 =\begin{pmatrix}0&-1\\1&-\gamma\end{pmatrix},
 \qquad \K q=p,\quad \K p=-q-\gamma p.
\end{equation}
Its two distinct eigenvalues are $-a$ and $-b$.  Finite-dimensional
diagonalization gives constants $0<c\leq C<\infty$ with
\[
 c e^{-at}\leq\|S(t)\|\leq C e^{-at}
 \qquad(t\geq0).
\]
Together with \eqref{eq:centered-norm}, this proves that the original sharp
gap is $a$.

Now use that $\Pj$ has range in $\HH_1$.  Since $\K$ preserves every chaos,
$\Qj$ is the identity and $\Qj\K\Qj=\K$ on $\HH_n$ for every $n\geq2$.
Within the first chaos, $\Qj\HH_1=\operatorname{span}\{p\}$ and
\eqref{eq:first} gives
\[
 \Qj\K\Qj p=-\gamma p.
\]
Consequently
\begin{equation}\label{eq:decomp}
 E=\operatorname{span}\{p\}\oplus\bigoplus_{n\geq2}\HH_n,
 \qquad
 T_{\Qj}(t)=e^{-\gamma t}I_{\operatorname{span}\{p\}}
 \oplus\bigoplus_{n\geq2}T(t)|_{\HH_n}.
\end{equation}
Polynomial chaoses form a core for the Gaussian Ornstein--Uhlenbeck
generator.  Moreover $\langle\K f,q\rangle=\langle f,\K^*q\rangle$ with
$\K^*q=-p$, so the rank-one compression term extends boundedly to $H$.
Thus the closure of the displayed block operator is precisely the
orthogonal generator, justifying \eqref{eq:decomp} at the semigroup level.

Applying Lemma \ref{lem:chaos} once more gives the exact norm formula
\begin{equation}\label{eq:qnorm}
 \|T_{\Qj}(t)|_E\|
 =\max\{e^{-\gamma t},\sup_{n\geq2}\|S(t)\|^n\}
 =\max\{e^{-\gamma t},\|S(t)\|^2\}.
\end{equation}
Since $\gamma=a+b>2a$, equations \eqref{eq:qnorm} and
\eqref{eq:centered-norm} imply upper and lower bounds of order $e^{-2at}$.
Hence the orthogonal sharp gap is $2a$.

For completeness, the same calculation identifies the point spectral sets:
the nonconstant eigenvalues of $\K$ are
\[
 \{-(ma+nb):m,n\in\mathbb N_0,\ m+n\geq1\},
\]
whereas after compression the first-chaos slow mode is gone and the retained
nonzero values are
\[
 \{-\gamma\}\cup
 \{-(ma+nb):m,n\in\mathbb N_0,\ m+n\geq2\}.
\]
The rightmost elements are respectively $-a$ and $-2a$.

Finally, \eqref{eq:first} gives $\Qj\K q=p$, and \eqref{eq:decomp} gives
$T_{\Qj}(t)p=e^{-\gamma t}p$.  Since $\langle q,p\rangle=0$ and
$\|q\|=1$,
\[
 \langle q,\K T_{\Qj}(t)p\rangle
 =e^{-\gamma t}\langle q,-q-\gamma p\rangle
 =-e^{-\gamma t}.
\]
This proves \eqref{eq:force-memory} and completes the theorem.
\qed

\section{Verification, novelty, and limitations}

The included checker constructs $\K$ on the normalized product-Hermite basis
through total degree ten.  For $\gamma\in\{2.1,2.5,3,5,10\}$ it verifies all
eigenvalues $-(ma+nb)$, the exact symmetric-tensor semigroup norms at three
times, the compressed first-chaos block, and \eqref{eq:force-memory}.  All
checks pass.  These computations are regression tests; Lemma
\ref{lem:chaos} and the direct-sum argument are the proof.

A bounded novelty audit searched the run's lightweight indexes and the web by
the exact paper id/title and close combinations of ``orthogonal dynamics'',
``harmonic/underdamped Langevin'', ``spectral gap'', and $\lambda_{\Qj}$.
It also inspected the relevant parts of arXiv:2503.20457, 2604.20453, and
2607.25855.  Those papers address semigroup well-posedness or infinite-rank
Zwanzig projections; their harmonic examples do not state this stochastic
rank-one gap-doubling theorem.  No exact duplicate was found as of 26 August
2026.  This is bounded evidence, not an exhaustive literature claim.

The theorem answers only the first open problem on source page 10, after
repairing its sign convention.  It does not treat the infinite-rank Zwanzig
projection, singular potentials, or the requested stochastic-analysis proof.
The restriction $\gamma>2$ avoids the critically damped Jordan case and keeps
the exact-rate statement free of polynomial prefactors.

\textbf{Human review focus.} Verify the Wick-exponential passage to symmetric
tensor powers, the core/closure identification behind \eqref{eq:decomp}, and
the intended interpretation of the source's reversed inequality.  Subject to
that convention repair, the answer is complete.

\begin{thebibliography}{9}
\bibitem{Zhu2021}
Y.~Zhu,
\emph{A simple hypocoercivity analysis for the effective Mori--Zwanzig equation},
arXiv:2109.01990, 2021.
\end{thebibliography}

\end{document}
